\section{Reproducibility, data and code availability}\label{sec:repro}

\paragraph{Repository layout.} The work is in the EVision project repository.
\begin{itemize}
\item \texttt{docs/thesis/}: the topic, the experiment protocol with its ``Changes'' log, the shock definition, and the results documents for the exploratory run, the first hold-out run, the robustness analyses, the post hoc additions (S6 and S7) and the extension.
\item \texttt{docs/thesis/results/}: small CSV result tables per run (exploratory runs at warm-up 48 and 96, hold-out runs at warm-up 48 and 96, the instant-feedback diagnostics, the degenerate post-2025-06 run, the robustness analyses, the extension runs with per-base-model analyses, and the clean-month accuracy of Chronos).
\item \texttt{scripts/thesis/}: Node scripts for the data download, the shock rule and the validity checks; the dated, hashed snapshots under \texttt{data/}; the experiment engine under \texttt{exp/} (data loading and chaining, forecasters including Chronos, interval methods, evaluation, calibration-error contrasts, power simulation, robustness analyses, command-line entry point and tests); the partial reproduction of DRACP; the explanation-layer harness under \texttt{llm\_eval/}.
\item \texttt{backend/fastapi\_service/}: the platform service for calibrated inflation intervals and the explanation service.
\item \texttt{paper/}: this paper, the bibliography, the figures and tables generated from the result files, and the build instructions.
\end{itemize}
Full per-row outputs stay under \texttt{scripts/thesis/results/} and are not tracked; the small result tables used in this paper are copied to \texttt{docs/thesis/results/}.

\paragraph{Environment.} Python 3.12 virtual environment under \texttt{scripts/thesis/.venv} with pandas, numpy, scipy, statsmodels, scikit-learn, LightGBM and the \texttt{dracp} package (version 0.2.0 in the reproduction check). The machine default was Python 3.14 with only numpy, and the package lists Python 3.10 to 3.12. The Chronos model is \texttt{amazon/chronos-bolt-tiny}, used zero-shot; the package version is not recorded in the documents read and is to be verified.

\paragraph{Tests.} The test suite of the engine has 81 tests; the modules cover data and chaining, the absence of look-ahead in the forecasters, the timing of delayed and instant feedback, the equivalence of the independent ACI implementation with the package at $h=1$, calibration-error contrasts, Holm, TOST, the power simulation, the hold-out loader, the robustness analyses, the Chronos forecaster with its leakage flag, and the reading of results files (the country code NA). The tests of the explanation layer (service and harness) number 18, and 31 together with the 13 platform-service tests.

\paragraph{Commands.} Commands are run from \texttt{scripts/thesis} with the interpreter of the virtual environment.
\begin{itemize}
\item Tests: \texttt{-m pytest exp/tests -q}.
\item Exploratory run: \texttt{-m exp.run full --start 2000-01 --end 2025-06 --out explore\_v1 --resume}.
\item Hold-out runs, all with \texttt{--holdout --start 2000-01 --end 2026-06}: run A with \texttt{--out holdout\_ncal48}; run B with \texttt{--n-cal 96 --out holdout\_ncal96}; run C with \texttt{--instant-feedback --horizons 12 --models lgbm --conformal-models lgbm --out diag\_instant\_h12\_holdout}.
\item Calibration-error contrasts: \texttt{-m exp.run ce --in results/<run> --out results/<run> --sample exploratory|holdout [--window-mde]}.
\item Robustness analyses: \texttt{-m exp.run robust --in results/holdout\_ncal48 --sample holdout --instant results/diag\_instant\_h12\_holdout}.
\end{itemize}
The command lines of the sensitivity run at warm-up 96 on the design countries, of the instant-feedback diagnostic on the design countries, and of the extension runs (second hold-out panel at both warm-ups, the re-run of the earlier 20 countries from 2010-01, the per-base-model calibration analysis and the clean-month accuracy script) were reconstructed from the command-line options of the entry point and the run names and are not recorded verbatim in the project documents; they are to be verified before submission. The explanation-layer harness has a command that runs a deterministic fake client for a harness check (not a result) and a command for a real run that refuses to run without an explicit spending flag and makes 936 calls; the real run was made on one model through an OpenAI-compatible proxy, and its raw answers and the v1 and v2 tables are in \texttt{docs/thesis/results/llm\_eval\_qwen/} and \texttt{docs/thesis/results/llm\_eval\_qwen\_v2/} of the repository (summary in \texttt{docs/thesis/results-llm-eval-v1.md}).

\paragraph{Paper assets.} The figures and tables of this paper are produced from the committed result files by \texttt{scripts/thesis/paper/build\_paper\_assets.py}, and the paper is compiled as described in \texttt{paper/README.md}. Numbers quoted in the running text were taken from the project manuscript and result documents and were not generated by that script.

\paragraph{Hashes and commits.} The data hashes are in Appendix~\ref{app:data} and the commit hashes of each protocol step in Appendix~\ref{app:changes}. The shock-rule document is dated 2026-11-18, and the first commit that contains it is the one recorded in Appendix~\ref{app:changes}. The exploratory run took 695 seconds.

\paragraph{Availability.} The code, the protocol, the data snapshots with hashes and the result tables are in the project repository. The underlying series are public IMF data. No private data are used.
